\chapter{Introduction and Mental Model}
\label{ch:intro}

\glance{ABC is a \emph{command shell} wrapped around \emph{graph data structures}
for Boolean logic. One global object (the \emph{frame}) holds the current design
in one of two forms --- the classic network \id{Abc\_Ntk\_t} and the modern
\id{Gia\_Man\_t}. Every command reads that design, rewrites it, and puts it back.
Learn the frame, the two graphs, and command dispatch first; everything else is
an algorithm plugged into that skeleton.}

\section{What ABC is}

ABC (``A System for Sequential Logic Synthesis and Formal Verification'') is a
logic-level EDA tool from UC Berkeley (Alan Mishchenko et al.). It takes a Boolean
circuit --- gates, LUTs, or an And-Inverter Graph (AIG) --- and makes it
\emph{smaller}, \emph{shallower}, or \emph{mapped to a technology}, and it can
\emph{prove} that two circuits are equivalent.

ABC operates \textbf{above} physical design. In a typical flow an RTL frontend
(e.g.\ Yosys) produces a netlist, ABC optimizes and maps it, and a physical tool
such as OpenROAD places and routes the result. Knowledge of ABC transfers
directly to the \cmd{abc} pass scripts used inside Yosys.

\begin{center}
\begin{tabular}{@{}lll@{}}
\toprule
\textbf{Problem} & \textbf{Typical command} & \textbf{Chapter} \\
\midrule
Too many gates / too deep      & \cmd{resyn2}, \cmd{\&syn2}, \cmd{rewrite}, \cmd{balance} & \ref{ch:opt} \\
Map to FPGA LUTs               & \cmd{if -K 6}, \cmd{\&if -K 6}                            & \ref{ch:map} \\
Map to standard cells          & \cmd{map}, \cmd{amap}                                     & \ref{ch:map} \\
Are two designs equal?         & \cmd{cec}, \cmd{\&cec}, \cmd{fraig}                       & \ref{ch:proof} \\
Sequential property proof      & \cmd{pdr}, BMC                                            & \ref{ch:proof} \\
Read / write a format          & \cmd{read\_*}, \cmd{write\_*}                             & \ref{ch:io} \\
\bottomrule
\end{tabular}
\end{center}

\section{The one picture that matters}

\begin{lstlisting}[style=abcshell]
                    +--------------------------------------+
                    |         Abc_Frame_t (global)         |
                    |  command table . aliases . libraries |
                    +------------------+-------------------+
                                       |
              +------------------------+------------------------+
              v                                                 v
       pNtkCur : Abc_Ntk_t                              pGia : Gia_Man_t
       "classic world"                                  "ABC9 / & world"
       rewrite, resyn2, if, map                         &syn2, &if, &b, &fraig
              |                                                 |
              +---------------------- &get / &put --------------+
                                       |
                                       v
                     read -> strash -> optimize -> map -> verify -> write
\end{lstlisting}

Two design representations live in the \emph{same process} and are toggled by the
bridge commands \cmd{\&get} (classic\,$\rightarrow$\,GIA) and \cmd{\&put}
(GIA\,$\rightarrow$\,classic). Table~\ref{tab:twoworlds} summarizes the split;
Chapters~\ref{ch:data} and~\ref{ch:aig} give the full detail.

\begin{table}[h]
\centering
\begin{tabularx}{\linewidth}{@{}lXX@{}}
\toprule
 & \textbf{Classic \id{Abc\_Ntk\_t}} & \textbf{GIA \id{Gia\_Man\_t}} \\
\midrule
Commands   & \cmd{rewrite}, \cmd{resyn2}, \cmd{if} & \cmd{\&b}, \cmd{\&syn2}, \cmd{\&if} \\
Node storage & object structs + ID vectors & one dense array, fanins as offsets \\
Scale        & up to \(\sim\)1M nodes & many-million-node designs \\
Held in frame & \id{pNtkCur} & \id{pGia} \\
Header       & \file{src/base/abc/abc.h} & \file{src/aig/gia/gia.h} \\
\bottomrule
\end{tabularx}
\caption{The two coexisting representations.}
\label{tab:twoworlds}
\end{table}

\section{The core idea behind almost every optimizer}

Nearly every optimization pass in ABC is a \textbf{greedy local rewrite} with
hand-tuned limits:

\begin{enumerate}
  \item Look at a \emph{local} piece of the graph (a node plus a small cut or cone).
  \item Try a \emph{functionally equivalent} replacement using fewer nodes or less depth.
  \item Accept it if the \emph{gain} is positive (greedy).
  \item Repeat over all nodes, usually inside a fixed recipe such as \cmd{resyn2}.
\end{enumerate}

Structural hashing keeps the AIG compact and lets rewrites \emph{reuse} existing
logic. Recipes like \cmd{resyn2} are not monolithic C algorithms --- they are
alias chains defined in \file{abc.rc} (Chapter~\ref{ch:recipes}).

\section{A first session}

\begin{lstlisting}[style=abcshell]
# batch mode: optimize and map i10 to 6-input LUTs
./abc -c "read i10.aig; strash; resyn2; if -K 6; print_stats; write_blif out.blif"

# the modern GIA path
./abc -c "read i10.aig; &get; &syn2; &if -K 6; &ps; &put; write_blif out.blif"
\end{lstlisting}

The metrics printed by \cmd{print\_stats} (alias \cmd{ps}) that you watch
constantly: \texttt{and} (AND-node count, an area proxy) and \texttt{lev} (logic
levels, a depth proxy).

\section{How to read this manual}

\begin{itemize}
  \item \textbf{Chapters \ref{ch:intro}--\ref{ch:data}} give the skeleton:
        architecture, command dispatch, and data structures. Read these first.
  \item \textbf{Chapters \ref{ch:aig}--\ref{ch:io}} cover the graph managers and
        how designs enter and leave ABC.
  \item \textbf{Chapters \ref{ch:opt}--\ref{ch:sat}} are the algorithms:
        optimization, mapping, verification, and the SAT engines they call.
  \item \textbf{Chapters \ref{ch:util}--\ref{ch:recipes}} cover the support
        libraries, scripts, and how to embed or extend ABC.
\end{itemize}

Throughout, \file{monospace blue} is a source path, \cmd{bold monospace} is a
shell command, \id{monospace} is a C identifier, and \pkg{teal} is a package
directory. All paths are relative to the ABC clone at \file{hardware\_eda/abc}.

\glance{Do \textbf{not} start reading in \pkg{sat/} or \pkg{bdd/} --- those are
engines ABC \emph{calls}, not the story. Always start from a command you can type,
then grep its name in \file{src/base/abci/abc.c} and follow the function pointer.}
